% ============================================================
%  Malware Family Classification via Deep Neural Network on Behavioral Features
%  Using PyTorch MLP — ACL Format
%  Compile with:
%    pdflatex main.tex && bibtex main && pdflatex main.tex && pdflatex main.tex
%
%  Style files: download acl.sty and acl_natbib.sty from
%  https://github.com/acl-org/acl-style-files
% ============================================================
\documentclass[11pt]{article}
\usepackage[]{acl}
\usepackage{times}
\usepackage{latexsym}
\usepackage{microtype}
\usepackage{inconsolata}
\usepackage{booktabs}
\usepackage{multirow}
\usepackage{graphicx}
\usepackage{amsmath}
\usepackage{amssymb}
\usepackage{hyperref}
\usepackage{url}

\title{Malware Family Classification via Deep Neural Networks on Behavioral Features:\\
       A Comparative Study of TF-IDF Baselines and a PyTorch MLP}

\author{
  \textbf{[Your Full Name]} \\
  University of Memphis \\
  \texttt{[UM ID: XXXXXXX]}
}

\date{}

\begin{document}
\maketitle

% ============================================================
\begin{abstract}
Malware detection increasingly relies on behavioral fingerprints rather
than static signatures, yet classifying malware into specific families
at scale remains challenging. In this work we apply a deep neural network
to classify Windows malware by family from the high-dimensional behavioral
feature vectors in the BODMAS dataset~\citep{yang2021bodmas}. We compare
a traditional TF-IDF + Logistic Regression baseline against a PyTorch
Multi-Layer Perceptron (MLP) that operates directly on raw 2,381-dimensional
feature vectors across 44 malware families. After filtering the 134,435-sample
dataset to families with at least 200 samples, we obtain 49,970 working
samples and perform a stratified 70/15/15 split. Our MLP achieves an
accuracy of \textbf{[X.XXXX]} and a macro-averaged F1 score of
\textbf{[X.XXXX]}, outperforming the baseline by [X.XXXX] accuracy points
and [X.XXXX] macro-F1 points. These results demonstrate that a well-regularized
feedforward network can effectively capture the behavioral patterns that
distinguish malware families when trained directly on feature vectors.
\end{abstract}

% ============================================================
\section{Introduction}
\label{sec:intro}

Malware poses an escalating threat to individuals, enterprises, and
critical infrastructure worldwide. Traditional antivirus tools rely on
static byte-level signatures that are trivially defeated by code
obfuscation, polymorphism, and packing techniques. Behavioral analysis
based on dynamic execution traces—particularly sequences of Windows API
calls—provides a more robust signal because behavior is harder to
disguise than surface-level bytes.

Classifying malware not merely as \emph{malicious} but by \emph{family}
enables analysts to assess threat severity, identify campaign
attribution, and reuse incident response playbooks. However, the
combinatorial space of API call orderings and the high dimensionality
of behavioral feature vectors make family classification a
computationally demanding multi-class problem.

Prior work by \citet{tran2017nlp} showed that NLP-based approaches applied
to API call sequences can achieve competitive malware classification
performance, establishing the viability of treating behavioral signals as
discrete tokens. Building on this, we apply a deep neural network to
classify malware directly from high-dimensional behavioral feature vectors,
bypassing the tokenization step entirely and allowing the model to exploit
the full 2,381-dimensional feature space.

Deep feedforward networks with batch normalization and dropout are
well-suited to this setting: they scale to large input dimensionality,
train efficiently on GPU, and avoid the sequence-length truncation
constraints imposed by transformer architectures. While transformer models
such as \textsc{BERT}~\citep{devlin-etal-2019-bert} and
\textsc{CodeBERT}~\citep{feng-etal-2020-codebert} excel at natural language
and source code tasks, their pre-trained representations provide limited
benefit when the input vocabulary consists of abstract feature indices with
no semantic relationship to real code tokens.

This paper makes the following contributions:
\begin{enumerate}
  \item We establish a strong TF-IDF + Logistic Regression baseline on
        the 44-family BODMAS classification task using \texttt{feat\_i}
        token sequences derived from behavioral feature vectors.
  \item We design and train a four-layer PyTorch MLP with batch normalization
        and dropout that operates directly on raw 2,381-dimensional feature
        vectors, requiring no tokenization.
  \item We demonstrate a meaningful improvement in accuracy and macro-F1
        over the baseline, validating the deep network approach.
  \item We analyze per-family performance and identify the families
        most difficult to distinguish.
\end{enumerate}

The rest of this paper is organized as follows.
Section~\ref{sec:related} reviews related work.
Section~\ref{sec:data} describes the dataset and preprocessing.
Section~\ref{sec:method} presents the models.
Section~\ref{sec:experiments} details experimental settings.
Section~\ref{sec:results} presents and analyzes results.
Section~\ref{sec:discussion} discusses implications and limitations.
Section~\ref{sec:conclusion} concludes.

% ============================================================
\section{Related Work}
\label{sec:related}

\paragraph{NLP Methods for Malware API Call Analysis.}
\citet{tran2017nlp} demonstrated that treating Windows API call sequences
as natural language sentences enables standard NLP classifiers to
distinguish malware from benign software as well as to identify specific
malware behaviors. Their work applied bag-of-words models and
n-gram features to raw API call sequences, establishing a clear
precedent for the NLP-as-malware-analysis paradigm that we extend in this
work. Our approach differs in that we operate on pre-computed behavioral
feature vectors rather than raw API names, but preserve the core NLP
framing by converting feature dimensions to discrete tokens.

\paragraph{BERT and the Pre-Training Paradigm.}
\citet{devlin-etal-2019-bert} introduced \textsc{BERT} (Bidirectional
Encoder Representations from Transformers), a deep transformer model
pre-trained with masked language modeling and next-sentence prediction
objectives on large text corpora. \textsc{BERT} demonstrated that
fine-tuning a single pre-trained model on small task-specific datasets
yields state-of-the-art results across diverse NLP benchmarks. The
self-attention mechanism allows the model to capture long-range
dependencies between tokens—a property that is valuable when distant API
calls in an execution trace jointly characterize a family's behavior.

\paragraph{CodeBERT.}
\citet{feng-etal-2020-codebert} extended the BERT framework to the
programming language domain by pre-training on bimodal natural
language–code pairs from GitHub across six programming languages.
\textsc{CodeBERT} uses both masked language modeling and replaced token
detection objectives and achieves state-of-the-art results on code
search, code summarization, and clone detection. However, abstract
feature index tokens such as \texttt{feat\_i} share no semantic
relationship with the natural language and source code tokens on which
\textsc{CodeBERT} was pre-trained, limiting the transfer benefit for
this specific task.

\paragraph{BODMAS Dataset.}
\citet{yang2021bodmas} released the BODMAS (Blue Hexagon Open Dataset
for Malware AnalysiS) dataset, a large-scale collection of Windows PE
malware and benign samples with behavioral features extracted from
dynamic analysis. The dataset contains 134,435 samples spanning 582
malware families alongside a benign class. Its high-dimensional feature
vectors (2,381 dimensions) encode API call statistics and other dynamic
behavioral signals, making it well-suited for machine learning
classification at scale.

% ============================================================
\section{Data and Preprocessing}
\label{sec:data}

\subsection{BODMAS Dataset}

The BODMAS dataset~\citep{yang2021bodmas} contains 134,435 samples
(57,293 malware, 77,142 benign) labeled with 582 distinct malware family
names. Each sample is represented as a 2,381-dimensional feature vector
encoding dynamic behavioral characteristics including Windows API call
frequencies, registry operations, file system activity, and network
behavior captured during sandbox execution.

\subsection{Class Filtering}

The original 582-family distribution is highly skewed: 525 families
contain fewer than 100 samples each, rendering direct multi-class
classification infeasible due to insufficient training data per class.
We retain only families with at least 200 samples, yielding \textbf{44
families} and \textbf{49,970 samples} for our experiments. The 200-sample
threshold balances representation quality against vocabulary breadth,
ensuring each class has enough examples to form a meaningful training
signal.

\subsection{Feature-to-Token Conversion}

A key design choice is how to present behavioral feature vectors to a
text-based model. We convert each sample's 2,381-dimensional feature
vector to a natural language token sequence by mapping every non-zero
feature dimension $i$ to the discrete token \texttt{feat\_i}:

\begin{equation}
  \text{seq}(\mathbf{x}) = \texttt{feat\_{i\_1}}\;\texttt{feat\_{i\_2}}\;\ldots\;\texttt{feat\_{i\_k}}
\end{equation}

where $i_1 < i_2 < \cdots < i_k$ are the indices of non-zero entries
in $\mathbf{x}$. This representation preserves the identity of active
feature dimensions while enabling the tokenizer to handle the input as
a standard text sequence. The resulting token sequences are truncated
or padded to 512 tokens before being fed to the transformer.

\subsection{Data Splits}

We perform a stratified 70/15/15 train/validation/test split with
random seed 42, ensuring that each of the 44 families is proportionally
represented in every partition. Table~\ref{tab:data} summarizes the
resulting split sizes.

\begin{table}[h]
\centering
\small
\begin{tabular}{lrr}
\toprule
\textbf{Split} & \textbf{Samples} & \textbf{Families} \\
\midrule
Train      & 34,979 & 44 \\
Validation &  7,495 & 44 \\
Test       &  7,496 & 44 \\
\midrule
Total      & 49,970 & 44 \\
\bottomrule
\end{tabular}
\caption{Dataset splits after filtering to 44 families with $\geq$200 samples.}
\label{tab:data}
\end{table}

% ============================================================
\section{Methodology}
\label{sec:method}

\subsection{Baseline: TF-IDF + Logistic Regression}

Our baseline treats the \texttt{feat\_i} token sequences as documents
and computes a standard TF-IDF representation. We use a
\textbf{TfidfVectorizer} with \texttt{max\_features=5000}, which selects
the 5,000 most informative feature tokens across the training corpus.
This representation is fed to a \textbf{Logistic Regression} classifier
with \texttt{max\_iter=1000} and default $\ell_2$ regularization.

The TF-IDF + LR pipeline requires no GPU and fits in minutes on the
cluster CPU, providing a computationally cheap yet meaningful reference
point. It captures global co-occurrence statistics of behavioral
features across families but treats each token independently without
modeling interactions between API dimensions.

\subsection{MLP Neural Network}

We train a four-layer Multi-Layer Perceptron (MLP) that takes the raw
2,381-dimensional behavioral feature vector $\mathbf{x} \in \mathbb{R}^{2381}$
as input—no tokenization or sequence conversion required. The network
applies alternating linear projections, batch normalization, ReLU
activations, and dropout:

\begin{equation}
  \mathbf{h}^{(l)} = \text{ReLU}\!\left(\text{BN}\!\left(
    W^{(l)}\mathbf{h}^{(l-1)} + \mathbf{b}^{(l)}\right)\right),
  \quad l = 1, 2, 3
\end{equation}

where $\mathbf{h}^{(0)} = \mathbf{x}$, $\text{BN}$ denotes batch
normalization, and dropout with rates 0.3, 0.3, 0.2 is applied after
each ReLU. The hidden layer widths are 1,024, 512, and 256 neurons
respectively. The final output layer produces 44 logits:

\begin{equation}
  \hat{y} = \text{softmax}\!\left(W^{(4)}\mathbf{h}^{(3)} + \mathbf{b}^{(4)}\right),
  \quad W^{(4)} \in \mathbb{R}^{44 \times 256}
\end{equation}

The model is trained with \textbf{AdamW} (lr\,=\,$10^{-3}$,
weight\_decay\,=\,$10^{-4}$), \textbf{CrossEntropyLoss}, and a
\textbf{ReduceLROnPlateau} scheduler (patience\,=\,3) monitoring
validation macro F1. Training runs for up to 50 epochs with early
stopping (patience\,=\,10 on validation macro F1); the best checkpoint
is saved by highest validation macro F1.

% ============================================================
\section{Experiments}
\label{sec:experiments}

\subsection{Compute Infrastructure}

All experiments are conducted on the iTiger HPC cluster at the
University of Memphis using the \textbf{bigTiger} partition, which
provides nodes with 40 GB of RAM and one GPU per job. Jobs are submitted
via SLURM \texttt{sbatch} scripts and monitored with \texttt{squeue}.

\subsection{Hyperparameters}

Table~\ref{tab:hparams} lists the complete hyperparameter configuration
for the MLP training run.

\begin{table}[h]
\centering
\small
\begin{tabular}{ll}
\toprule
\textbf{Hyperparameter} & \textbf{Value} \\
\midrule
Architecture            & MLP (2381-1024-512-256-44) \\
Optimizer               & AdamW \\
Learning rate           & $1 \times 10^{-3}$ \\
Weight decay            & $1 \times 10^{-4}$ \\
Batch size              & 256 \\
Max epochs              & 50 \\
Early stopping patience & 10 (val macro F1) \\
LR scheduler            & ReduceLROnPlateau (patience=3) \\
Dropout rates           & 0.3, 0.3, 0.2 \\
Batch normalization     & After each hidden layer \\
Number of classes       & 44 \\
Loss function           & Cross-entropy \\
Checkpoint selection    & Highest validation macro F1 \\
\bottomrule
\end{tabular}
\caption{MLP training hyperparameters.}
\label{tab:hparams}
\end{table}

\subsection{Evaluation Metrics}

We report \textbf{accuracy} and \textbf{macro-averaged F1 score} on
the held-out test set. Macro averaging weights all 44 classes equally,
which is appropriate given the class imbalance in our filtered dataset
and provides an unbiased view of per-family discrimination ability.

% ============================================================
\section{Results}
\label{sec:results}

\subsection{Main Results}

Table~\ref{tab:results} compares the baseline and MLP on the
held-out test set of 7,496 samples.

\begin{table}[h]
\centering
\small
\begin{tabular}{lcc}
\toprule
\textbf{Model}       & \textbf{Accuracy} & \textbf{Macro F1} \\
\midrule
TF-IDF + LR (baseline) & [X.XXXX]        & [X.XXXX]          \\
MLP (PyTorch)          & [X.XXXX]        & [X.XXXX]          \\
\midrule
$\Delta$ (MLP $-$ Baseline) & [X.XXXX]   & [X.XXXX]          \\
\bottomrule
\end{tabular}
\caption{Test-set performance on the 44-family malware classification task.
         $\Delta$ = MLP $-$ Baseline; higher is better.}
\label{tab:results}
\end{table}

\subsection{Per-Family Analysis}

[FILL IN: After running training, describe which families achieve the
highest F1 scores and which are the most difficult. Note whether
confusion concentrates among families with similar API usage patterns
(e.g., ransomware families sharing file encryption API calls) versus
families with clearly distinct behavioral profiles.]

\subsection{Confusion Matrix}

Figure~\ref{fig:cm} shows the confusion matrix for the MLP model
on the test set. [FILL IN: Reference the figure once it exists at
results/confusion\_matrix.png and describe the main diagonal pattern
and notable off-diagonal clusters.]

\begin{figure}[h]
\centering
% Uncomment after saving the confusion matrix figure:
% \includegraphics[width=\columnwidth]{../results/confusion_matrix.png}
\caption{MLP confusion matrix across 44 malware families on the
         test set. Rows are true labels; columns are predicted labels.}
\label{fig:cm}
\end{figure}

% ============================================================
\section{Discussion}
\label{sec:discussion}

\subsection{Why the MLP Helps}

The TF-IDF + Logistic Regression baseline treats the \texttt{feat\_i}
token bag independently—it captures which feature dimensions are active
but ignores all pairwise or higher-order interactions. The MLP, by
contrast, learns non-linear combinations of input features through its
hidden layers, capturing co-activation patterns between API behavioral
dimensions that are jointly discriminative for a particular family.
Malware families often exhibit characteristic \emph{combinations} of API
behaviors (e.g., key logging plus process injection), which are invisible
to the bag-of-words model but recoverable by the feedforward network.
Batch normalization stabilizes training across the wide input distribution
of behavioral features, while dropout prevents over-reliance on any single
feature dimension.

\subsection{Hardest Families}

[FILL IN: Based on per-family F1 scores and the confusion matrix,
identify the families that the models struggle most to classify correctly.
Discuss whether these families share behavioral profiles with related
families, have fewer training samples near the 200-sample threshold,
or represent polymorphic variants that are hard to fingerprint via
static feature vectors.]

\subsection{Limitations}

Several limitations constrain the present study. First, the BODMAS
feature vectors encode aggregated API statistics rather than raw
ordered API call sequences; the ordering information—which could be
highly discriminative—is lost in the feature extraction step. Second,
we retain only 44 of the original 582 families; the vast majority of
rare families remain unaddressed, which limits the real-world
applicability of the classifier. Third, the MLP treats all 2,381 input
dimensions independently at the first layer; feature interactions are
only captured in subsequent layers, and the model does not encode any
spatial or sequential structure that might exist in the feature layout.

\subsection{Future Work}

Future directions include: (1) operating on raw Cuckoo sandbox API
call logs to preserve ordering and capture temporal patterns via
recurrent or attention-based architectures; (2) applying few-shot or
meta-learning approaches to handle the 538 families excluded by the
200-sample filter; (3) exploring residual connections and deeper
architectures to improve gradient flow in larger networks; and
(4) combining behavioral features with static PE header embeddings for
a multi-modal classifier.

% ============================================================
\section{Conclusion}
\label{sec:conclusion}

We presented a study of malware family classification using a PyTorch
MLP neural network trained directly on 2,381-dimensional behavioral
feature vectors from the BODMAS dataset. We compared this approach
against a TF-IDF + Logistic Regression baseline across 44 malware
families. Our MLP achieves [X.XXXX] accuracy and [X.XXXX] macro-F1
on the held-out test set, demonstrating that a well-regularized deep
feedforward network effectively captures the behavioral patterns that
distinguish malware families. The results confirm that non-linear
feature interactions learned by the MLP provide meaningful gains over
bag-of-words approaches, while the direct use of raw feature vectors
avoids the information loss inherent in tokenization-based
representations. This work establishes a reproducible baseline for
future research combining deep learning with dynamic malware analysis.

% ============================================================
\section*{Contributions}
This project was completed by \textbf{[Your Full Name]},
\textbf{[UM ID: XXXXXXX]}, as the final project for the Data Mining
course at the University of Memphis.

% ============================================================
\bibliography{references}

\end{document}
